% Template: one section of a chapter.
% Copy to documentation/tex/notes/<topic>/sections/<section>.tex.
% This file contains prose only: no preamble, and no \section command --
% the sectioning lives in the parent chapter file.

Open directly with the prose.
Break the source at clause boundaries,
one sentence or one clause per line,
so that a diff shows exactly which sentence changed.

Define terms with \emph{} on first use,
and give every symbol a macro in include/notation.tex
rather than writing it out here.

\begin{defi}
\label{def.something}
The \emph{something} at time $t$ is
\begin{equation}
\label{eq.something}
x_t := 0.
\end{equation}
\end{defi}

Refer back with the noun spelled out:
Definition \ref{def.something} and equation \eqref{eq.something},
and cite as \cite{Bou18tra}.

\begin{remark}
\label{remark.something}
Use remarks for the caveat that a careful reader would raise,
and that a student would otherwise get wrong in code.
\end{remark}

\subsection{A subsection, if the section needs one}
\label{sec.aSubsection}

Subsections do live in the section file;
only \section belongs to the parent.
